% Generated draft: requires mathematical and source review.
\documentclass[11pt]{article}
\usepackage[margin=25mm]{geometry}
\usepackage{fontspec}
\setmainfont{Latin Modern Roman}
\usepackage{amsmath,amssymb,mathtools,unicode-math}
\setmathfont{Latin Modern Math}
\usepackage{enumitem,needspace,xurl,hyperref}
\setlength{\parindent}{0pt}
\setlength{\parskip}{5pt}
\setlength{\emergencystretch}{3em}
\newcommand{\N}{\mathbb{N}}
\newcommand{\Z}{\mathbb{Z}}
\newcommand{\Q}{\mathbb{Q}}
\newcommand{\R}{\mathbb{R}}
\newcommand{\C}{\mathbb{C}}
\newcommand{\F}{\mathbb{F}}
\newcommand{\A}{\mathbb{A}}
\newcommand{\PP}{\mathbb{P}}
\newcommand{\E}{\mathbb{E}}
\newcommand{\eps}{\varepsilon}
\newcommand{\dd}{\,\mathrm{d}}
\begin{document}
\section*{UnsolvedMath ID: 30001371}
\textbf{Source problem code:} catalogue entry
\subsection{Definitions and mathematical statement}
For terminology, let \(R\) be a ring with identity and let \((C_\bullet,d_\bullet)\) be a chain complex of left \(R\)-modules, with maps \(d_n:C_n\to C_{n-1}\) satisfying \(d_n\circ d_{n+1}=0\) for every \(n\in\Z\). Its cycles, boundaries, and homology in degree \(n\) are
\[
Z_n(C_\bullet)=\ker(d_n),\qquad B_n(C_\bullet)=\operatorname{im}(d_{n+1}),\qquad H_n(C_\bullet)=Z_n(C_\bullet)/B_n(C_\bullet).
\]
This module setting supplies standard terminology only; the source does not specify an ambient category or restrict its questions to this setting.

\textbf{Research questions (as supplied).} Determine when homology can be computed constructively, what algebraic conditions are required for such computation, and which results and techniques of homological algebra admit constructive formulations and implementations in software or proof assistants.

\textbf{Unresolved formulation qualifications.} The supplied passage describes a research programme rather than a single quantified conjecture. It does not specify the algebraic objects, their effective presentations, the degrees or boundedness of complexes, the permitted computational operations, or the required representation of the output homology. It also leaves unspecified the meaning of constructive computation, the intended homological results and techniques, and the software or foundational systems in which implementations are sought. Accordingly, no universal algorithmic assertion, necessary-and-sufficient criterion, or specific formalization theorem is asserted here.
\subsection{Short English statement}
The problem asks for algebraic settings in which homology can be computed constructively, for the hypotheses enabling those computations, and for constructive implementations of homological methods in software or proof assistants. The supplied formulation is a broad research programme whose precise computational and foundational targets remain unspecified.
\subsection{Sources}
\begin{itemize}
\item Local registry: \texttt{data/unsolvedmath/problems.json}, with \texttt{catalog/problems.json} as the supplementary registry; permanent UnsolvedMath ID \(30001371\), source code \texttt{catalogue entry}, title \emph{Constructive Homological Algebra}. This draft uses the supplied complete source record; registry files were not independently inspected.
\item Supplied report citation: \emph{Mini-Workshop: Formal Methods in Commutative Algebra: A View Toward Constructive Homological Algebra} (2010), \emph{Oberwolfach Reports}, pp.\ 2731--2770, DOI \texttt{10.4171/owr/2009/50}. Supplied source URL: \url{https://doi.org/10.4171/owr/2009/50}. The background describes these reports as documenting mathematical workshops and their open-problem discussions. No individual proposer is supplied.
\item Supplied literature source, labelled ``original Oberwolfach report 4133'': \url{https://doi.org/10.4171/owr/2009/50}.
\item Literature assessment dated exactly \texttt{2026-08-21}: the supplied curation labels the problem open and describes the passage as broad questions about computability, algebraic conditions, implementations, and proof assistants, rather than a resolved single theorem. This is supplied dated curation, not an independently verified literature assessment. The record requests independent verification of time-sensitive status and cited claims.
\item The record labels the statement exact and reports a manual review finding the extracted statement self-contained, mathematically coherent, and free of LaTeX or report-layout artifacts. These are supplied verification claims, not independent verification by this draft. The proposed year and year tag are 2009, whereas the supplied report citation gives 2010; the distinction is unresolved here.
\end{itemize}
\end{document}
